\makeabstract
    {
    Synthesis and Evaluation of Novel Alkyl-Sulfonated Near Infrared-II Fluorophores for Cancer Imaging
    }
    {
    Divine Kengne\textsuperscript{1,2}, Alex Vanover\textsuperscript{1}, Kayla Barekat\textsuperscript{1}, Ritesh Isuri\textsuperscript{1},  Anatoliy Popov\textsuperscript{1}, Jim Delikatny\textsuperscript{1}
    }
    {
    \textsuperscript{1} Department of Radiology, Perelman School of Medicine, University of Pennsylvania, Philadelphia, PA\\
\textsuperscript{2} Amherst College, Amherst, MA
    }
    {
    Fluorescence-Guided Surgery (FGS) has emerged as a sophisticated method of intraoperative molecular imaging that addresses the fundamental limitations of conventional intraoperative assessment (white-light visualization, palpation, and mental {\textquotedblleft}registration{\textquotedblright} of preoperative imaging) [1]. In FGS, a fluorescent probe is administered systemically prior to surgery and allowed to circulate and accumulate in the tumor over time.  The currently available fluorophores in FGS emit and absorb in the NIR-I region, which is characterized by shallow penetration and high autofluorescence. To address these limitations, our lab has developed a class of low-molecular-weight organic fluorophores, 4,4'-quinocyanines (QuCy), that extend emission into the NIR-II region [2]. Despite QuCy's success in achieving high-resolution visualization in murine models, its fluorescence is quenched when dissolved in aqueous media due to the formation of aggregates. For this reason, liposomal encapsulations were made to deliver this fluorophore to mice. To resolve this complication, we hypothesize that adding sulfonate groups to the QuCy head will increase their water solubility and prevent aggregation, thereby eliminating the need for a liposome for delivery. Preliminary data show that the alkyl-sulphonated QuCy (sbcQuCy) has a strong absorbance in NIR-II when dissolved in polar solvents like DMSO and Methanol. 
    }
    {
    With the synthesized dye, we carried out a cellular uptake assay, cytotoxicity assay, and intravenous injection in naive and tumor mice.
    }
    {
    Evaluation of sbcQuCy revealed efficient tumor uptake, low cytotoxicity, and high-resolution imaging
    }
    {
    With these findings, our hope is that sbcQuCy will serve as a clinically translatable NIR-II contrast agent
    }

    \newpage
    